\documentclass[10pt,a4paper]{amsart}
\usepackage{amsmath,amssymb,amsthm,graphicx}
\usepackage{verbatim}
\newenvironment{reference}{\comment}{\endcomment}
%\newenvironment{reference}{}{}
\newenvironment{slogan}{\comment}{\endcomment}
\newenvironment{history}{\comment}{\endcomment}

% For commutative diagrams we use Xy-pic
\usepackage[all]{xy}

% We use 2cell for 2-commutative diagrams.
\xyoption{2cell}
% Theorem environments.
%
\theoremstyle{plain}
\newtheorem{theorem}[subsection]{Theorem}
\newtheorem{proposition}[subsection]{Proposition}
\newtheorem{lemma}[subsection]{Lemma}

\theoremstyle{definition}
\newtheorem{definition}[subsection]{Definition}
\newtheorem{example}[subsection]{Example}
\newtheorem{exercise}[subsection]{Exercise}
\newtheorem{situation}[subsection]{Situation}

\theoremstyle{remark}
\newtheorem{remark}[subsection]{Remark}
\newtheorem{remarks}[subsection]{Remarks}

\numberwithin{equation}{subsection}

% Macros
%
\def\lim{\mathop{\mathrm{lim}}\nolimits}
\def\colim{\mathop{\mathrm{colim}}\nolimits}
\def\Spec{\mathop{\mathrm{Spec}}}
\def\Hom{\mathop{\mathrm{Hom}}\nolimits}
\def\Ext{\mathop{\mathrm{Ext}}\nolimits}
\def\SheafHom{\mathop{\mathcal{H}\!\mathit{om}}\nolimits}
\def\SheafExt{\mathop{\mathcal{E}\!\mathit{xt}}\nolimits}
\def\Sch{\mathit{Sch}}
\def\Mor{\mathop{\mathrm{Mor}}\nolimits}
\def\Ob{\mathop{\mathrm{Ob}}\nolimits}
\def\Sh{\mathop{\mathit{Sh}}\nolimits}
\def\NL{\mathop{N\!L}\nolimits}
\def\CH{\mathop{\mathrm{CH}}\nolimits}
\def\proetale{{pro\text{-}\acute{e}tale}}
\def\etale{{\acute{e}tale}}
\def\QCoh{\mathit{QCoh}}
\def\Ker{\mathop{\mathrm{Ker}}}
\def\Im{\mathop{\mathrm{Im}}}
\def\Coker{\mathop{\mathrm{Coker}}}
\def\Coim{\mathop{\mathrm{Coim}}}

% Boxtimes
%
\DeclareMathSymbol{\boxtimes}{\mathbin}{AMSa}{"02}

%
% Macros for moduli stacks/spaces
%
\def\QCohstack{\mathcal{QC}\!\mathit{oh}}
\def\Cohstack{\mathcal{C}\!\mathit{oh}}
\def\Spacesstack{\mathcal{S}\!\mathit{paces}}
\def\Quotfunctor{\mathrm{Quot}}
\def\Hilbfunctor{\mathrm{Hilb}}
\def\Curvesstack{\mathcal{C}\!\mathit{urves}}
\def\Polarizedstack{\mathcal{P}\!\mathit{olarized}}
\def\Complexesstack{\mathcal{C}\!\mathit{omplexes}}
% \Pic is the operator that assigns to X its picard group, usage \Pic(X)
% \Picardstack_{X/B} denotes the Picard stack of X over B
% \Picardfunctor_{X/B} denotes the Picard functor of X over B
\def\Pic{\mathop{\mathrm{Pic}}\nolimits}
\def\Picardstack{\mathcal{P}\!\mathit{ic}}
\def\Picardfunctor{\mathrm{Pic}}
\def\Deformationcategory{\mathcal{D}\!\mathit{ef}}
\usepackage{cleveref}
\crefname{originalreference}{original reference}{original references}
\Crefname{originalreference}{Original reference}{Original references}
\makeatletter
\newlabel{lemma-ok-for-projective-bundle}{{0FIN}{source}{}{sourcepointer-1}{}}
\newlabel{lemma-ok-for-projective-bundle@cref}{{[originalreference][0][]0FIN}{[1][0][]source}{}{}{}}
\makeatother
\title{Compatibility with pushforward along the projection from a Grassmannian}
\author{The Stacks Project Authors}
\date{stacks-project:weil.tex:local-snapshot-20261012; GNU FDL 1.2 or later}
\begin{document}
\maketitle
\noindent\textit{Editorial note.} All delivered mathematical text is the author's own, copied from the pinned Stacks Project chapter file weil.tex; no statement or proof step was rewritten, reordered or removed.\par\medskip
\noindent\textit{Editorial note.} The counted proof is the author's complete argument: the extreme cases as projective space bundles, the reduction through the partial flag variety, and the pushforward argument concluding the general case.\par\medskip
\noindent\textit{Editorial note.} The label lemma-grassmanian names the counted statement, and the proof environment immediately following it in the source is delivered verbatim as its proof.\par\medskip
\noindent\textit{Editorial note.} The counted statement is carried by the author's own lemma environment in the chapter source.\par\medskip
\noindent\textit{Editorial note.} Pre-open gate: python operations/agent-authoring/tools/precard\_gate.py --source weil.tex --label lemma-grassmanian --start 5227 --end 5255 --local-id 28154 returned RESULT=OK (L28154 free, occupancy FREE, public interval clear). Pre-assembly tag self-check: canonical tag resolved and not present in reference/public-items.jsonl.\par\medskip
\noindent\textit{Editorial note.} The statement and proof are copied from the pinned chapter file; the cited inputs are kept as external pointers with their canonical Stacks tags.\par\medskip
\noindent\textit{Allowed standard prerequisites (editorial).}\begin{itemize}
\item Compatibility with pushforward along a projective space bundle. Same chapter, the lemma that the cycle class map commutes with pushforward along the projectivization of a vector bundle. Provides the case of a projective space bundle, which is the base case of the reduction.
\item Partial flag varieties as iterated projective space bundles. The standard description of the partial flag variety as an iterated projective space bundle over the Grassmannian and over the base field. Produces the auxiliary scheme over which both projections become compositions of projective space bundles.
\end{itemize}\medskip
\ifdefined\hypertarget\hypertarget{sourcepointer-1}{}\fi
\noindent\textit{Source reference.} Original reference 0FIN: sources/stacks-weil/pinned-original-source.tex, source lines 5026--5026. Prerequisite pointer invoked at proof line(s) 5238; it names a standard fact of the same chapter that the counted argument uses, and every step of the counted proof built on it is retained verbatim in the delivered body.\par\medskip
\begin{lemma}
\label{lemma-grassmanian}
Assume given data (D0), (D1), and (D2') satisfying axioms (A1) -- (A7).
Given integers $0 < l < n$ and a nonempty equidimensional
smooth projective scheme $X$ over $k$ consider the projection morphism
$p : X \times \mathbf{G}(l, n) \to X$.
Then $\gamma$ commutes with pushforward along $p$.
\end{lemma}

\begin{proof}
If $l = 1$ or $l = n - 1$ then $p$ is a projective bundle and
the result follows from Lemma \ref{lemma-ok-for-projective-bundle}.
In general there exists a morphism
$$
h : Y \to X \times \mathbf{G}(l, n)
$$
such that both $h$ and $p \circ h$ are compositions of projective
space bundles. Namely, denote $\mathbf{G}(1, 2, \ldots, l; n)$
the partial flag variety. Then the morphism
$$
\mathbf{G}(1, 2, \ldots, l; n) \to \mathbf{G}(l, n)
$$
is a compostion of projective space bundles and similarly the
structure morphism $\mathbf{G}(1, 2, \ldots, l; n) \to \Spec(k)$
is of this form. Thus we may set $Y = X \times \mathbf{G}(1, 2, \ldots, l; n)$.
Since every cycle on $X \times \mathbf{G}(l, n)$ is the pushforward of
a cycle on $Y$, the result for $Y \to X$ and the result for
$Y \to X \times \mathbf{G}(l, n)$ imply the result for $p$.
\end{proof}

\end{document}
